% Options for packages loaded elsewhere
\PassOptionsToPackage{unicode}{hyperref}
\PassOptionsToPackage{hyphens}{url}
%
\documentclass[
]{book}
\usepackage{amsmath,amssymb}
\usepackage{lmodern}
\usepackage{iftex}
\ifPDFTeX
  \usepackage[T1]{fontenc}
  \usepackage[utf8]{inputenc}
  \usepackage{textcomp} % provide euro and other symbols
\else % if luatex or xetex
  \usepackage{unicode-math}
  \defaultfontfeatures{Scale=MatchLowercase}
  \defaultfontfeatures[\rmfamily]{Ligatures=TeX,Scale=1}
  \setmainfont[]{Libre Baskerville}
\fi
% Use upquote if available, for straight quotes in verbatim environments
\IfFileExists{upquote.sty}{\usepackage{upquote}}{}
\IfFileExists{microtype.sty}{% use microtype if available
  \usepackage[]{microtype}
  \UseMicrotypeSet[protrusion]{basicmath} % disable protrusion for tt fonts
}{}
\makeatletter
\@ifundefined{KOMAClassName}{% if non-KOMA class
  \IfFileExists{parskip.sty}{%
    \usepackage{parskip}
  }{% else
    \setlength{\parindent}{0pt}
    \setlength{\parskip}{6pt plus 2pt minus 1pt}}
}{% if KOMA class
  \KOMAoptions{parskip=half}}
\makeatother
\usepackage{xcolor}
\usepackage[left=4cm, right=3cm, top=2.5cm, bottom=2.5cm]{geometry}
\usepackage{listings}
\newcommand{\passthrough}[1]{#1}
\lstset{defaultdialect=[5.3]Lua}
\lstset{defaultdialect=[x86masm]Assembler}
\usepackage{longtable,booktabs,array}
\usepackage{calc} % for calculating minipage widths
% Correct order of tables after \paragraph or \subparagraph
\usepackage{etoolbox}
\makeatletter
\patchcmd\longtable{\par}{\if@noskipsec\mbox{}\fi\par}{}{}
\makeatother
% Allow footnotes in longtable head/foot
\IfFileExists{footnotehyper.sty}{\usepackage{footnotehyper}}{\usepackage{footnote}}
\makesavenoteenv{longtable}
\usepackage{graphicx}
\makeatletter
\def\maxwidth{\ifdim\Gin@nat@width>\linewidth\linewidth\else\Gin@nat@width\fi}
\def\maxheight{\ifdim\Gin@nat@height>\textheight\textheight\else\Gin@nat@height\fi}
\makeatother
% Scale images if necessary, so that they will not overflow the page
% margins by default, and it is still possible to overwrite the defaults
% using explicit options in \includegraphics[width, height, ...]{}
\setkeys{Gin}{width=\maxwidth,height=\maxheight,keepaspectratio}
% Set default figure placement to htbp
\makeatletter
\def\fps@figure{htbp}
\makeatother
\setlength{\emergencystretch}{3em} % prevent overfull lines
\providecommand{\tightlist}{%
  \setlength{\itemsep}{0pt}\setlength{\parskip}{0pt}}
\setcounter{secnumdepth}{5}
\usepackage{float}
\floatplacement{figure}{H}
\usepackage{booktabs}
\usepackage{setspace}
\usepackage{fontspec}
\onehalfspacing
\usepackage{xcolor}
\usepackage{titlesec}
\usepackage{xfrac}
% Make caption labels bold: https://bookdown.org/yihui/rmarkdown-cookbook/latex-extra.html
\usepackage[labelfont={bf}]{caption} 
% force floats forward: https://bookdown.org/yihui/rmarkdown-cookbook/figure-placement.html
\usepackage{flafter}
% https://tex.stackexchange.com/questions/279/how-do-i-ensure-that-figures-appear-in-the-section-theyre-associated-with
\usepackage{placeins}
\lstset{
  breaklines=true
}
\usepackage{wrapfig}
\usepackage{lipsum}
\usepackage{caption}
\captionsetup[figure]{font=small}
% following code copied from https://bookdown.org/yihui/rmarkdown-cookbook/figure-placement.html#fnref10 
\renewcommand{\topfraction}{.85}
\renewcommand{\bottomfraction}{.7}
\renewcommand{\textfraction}{.15}
\renewcommand{\floatpagefraction}{.66}
\setcounter{topnumber}{3}
\setcounter{bottomnumber}{3}
\setcounter{totalnumber}{4}

%% https://stackoverflow.com/questions/3275770/modifying-section-to-make-it-colorful-with-latex
%\usepackage{color}
%\titleformat{\section}
%{\color{red}}


% https://tex.stackexchange.com/questions/10320/section-and-subsection-colors-using-titlesec
%\makeatletter
%\newcommand*\@secondofsix[6]{#2}
%\newcommand{\addtotitleformat}{%
%  \@ifstar{\addtotitleformat@star}{\addtotitleformat@nostar}}
%\newcommand\addtotitleformat@nostar[2]{%
%  \PackageError{titlesec}{non starred form of \string\addtotitleformat\space not supported}{}}
%\newcommand\addtotitleformat@star[2]{%
%  \expandafter\expandafter\expandafter\expandafter
%  \expandafter\expandafter\expandafter\def
%  \expandafter\expandafter\expandafter\expandafter
%  \expandafter\expandafter\expandafter\@currentsection@font
%  \expandafter\expandafter\expandafter\expandafter
%  \expandafter\expandafter\expandafter{%
%    \expandafter\expandafter\expandafter\@secondofsix
%       \csname ttlf@\expandafter\@gobble\string#1\endcsname}%
%  \titleformat*{#1}{\@currentsection@font#2}%
%}
%\makeatother
%
%\addtotitleformat*{\section}{\Huge\color{red}}
%\addtotitleformat*{\subsection}{\sffamily\color{blue}}
\usepackage{titling}
\usepackage{subcaption}
\pretitle{\begin{center} \includegraphics[width=2in,height=2in]{/Users/brettell/Documents/Repositories/PhD-thesis/book/figs/title/Arms_PembrokeCollege_Cambridge.pdf}\LARGE\\ \bigskip \bigskip \bigskip }
\posttitle{\end{center}}
\predate{\begin{center} Pembroke College \\ \bigskip }
\postdate{\end{center} \begin{center} \bigskip \bigskip \bigskip This thesis is submitted for the degree of Doctor of Philosophy. \end{center} \centering \begin{figure}[!tpb] \centering \begin{subfigure}{0.49 \linewidth} \centering \vfill \includegraphics[height=0.45in]{/Users/brettell/Documents/Repositories/PhD-thesis/book/figs/title/cambridge_university2.pdf} \end{subfigure} \hfill \begin{subfigure}{0.49 \linewidth} \centering \includegraphics[height=0.55in]{/Users/brettell/Documents/Repositories/PhD-thesis/book/figs/title/EMBL_EBI_Logo_black.pdf} \end{subfigure} \end{figure} }
\usepackage{float}
\ifLuaTeX
  \usepackage{selnolig}  % disable illegal ligatures
\fi
\IfFileExists{bookmark.sty}{\usepackage{bookmark}}{\usepackage{hyperref}}
\IfFileExists{xurl.sty}{\usepackage{xurl}}{} % add URL line breaks if available
\urlstyle{same} % disable monospaced font for URLs
\hypersetup{
  pdftitle={Japanese courage: a genetic analysis of complex traits in medaka fish and humans},
  pdfauthor={Ian Narain Brettell},
  hidelinks,
  pdfcreator={LaTeX via pandoc}}

\title{Japanese courage: a genetic analysis of complex traits in medaka fish and humans}
\author{Ian Narain Brettell}
\date{30 September 2022}

\begin{document}
\maketitle

{
\setcounter{tocdepth}{1}
\tableofcontents
}
\hypertarget{acknowledgements}{%
\chapter*{Acknowledgements}\label{acknowledgements}}
\addcontentsline{toc}{chapter}{Acknowledgements}

Thanks.

\hypertarget{preface}{%
\chapter*{Preface}\label{preface}}
\addcontentsline{toc}{chapter}{Preface}

This thesis is the result of my own work and includes nothing which is the outcome of work done in collaboration except as declared in the preface and specified in the text.

It is not substantially the same as any work that has already been submitted before for any degree or other qualification except as declared in the preface and specified in the text.

It does not exceed the prescribed word limit for the Biology Degree Committee.

\hypertarget{abstract}{%
\chapter*{Abstract}\label{abstract}}
\addcontentsline{toc}{chapter}{Abstract}

\textbf{Japanese courage: a genetic analysis of complex traits in medaka fish and humans}

This thesis primarily explores how an individual's genes interact with the genes of their social companions to create differences in behaviour, using the Japanese medaka fish as a model organism. Chapter 1 sets out the introduction to the diverse topics covered in this thesis.

Chapter 2 describes several genomic characteristics of the Medaka Inbred Kiyosu-Karlsruhe (MIKK) panel, which comprises 80 inbred lines of medaka that were bred from a wild population residing in Kiyosu, southern Japan. In this chapter I plot the inbreeding trajectory of the MIKK panel, and analyse its evolutionary relationship with other previously established inbred medaka strains; degree of homozygosity; rate of linkage disequilibrium decay; repeat content; and structural variation, all which relate to its utility for the genetic mapping of complex traits.

In Chapter 3, I use a custom behavioural assay to characterise and classify bold-shy behaviours in 5 previously established inbred medaka lines. Here I describe the assay, assess its robustness against confounding factors, and apply a hidden markov model (HMM) to classify the fishes' behaviours across a spectrum of boldness-shyness based on their distance and angle of travel between time points. I describe how the different lines differ in their behaviours over the course of the assay (a direct genetic effect) and how the behaviour of a single ``reference'' line (\emph{iCab}) differs in the presence of different lines (a social genetic effect).

In Chapter 4, I explain how I applied this behavioural assay to the MIKK panel in order to identify lines that diverge in both their own bold-shy behaviours (the direct genetic effect) and the extent to which they transmit those behaviours onto their tank partners (the social genetic effect). I then describe how we used those divergent lines as the parental lines in a multi-way F2 cross in an attempt to isolate the genetic variants that are associated with both direct and social genetic effects.

In Chapter 5 I describe the bioinformatic processes and genetic association models used to map the variants associated with differences in the period of somite development, based on a separate F2 cross between the southern Japanese \emph{iCab} strain, and the northern Japanese \emph{Kaga} strain.

Finally, in Chapter 6, I compare and rank all complex traits in the GWAS Catalog based on the extent to which their associated alleles vary across global human populations, using the Fixation Index (Fst) as a metric and the 1000 Genomes dataset as a sample of global genetic variation. In this chapter I set out the bioinformatic pipelines used to process the data, present the distributions of Fst for trait-associated alleles across the genome, and use the Kolmogorov-Smirnov test to compare the distributions of Fst across different traits.

Altogether, this thesis describes some of the genomic characteristics of both medaka fish and humans, and how those variations relate to differences in complex traits, with a particular focus on the genetic causes of adaptive behaviours and the transmission of those behaviours onto one's social companions.

\end{document}
